\section{Limpieza y preparación}

\subsection{Código aplicado}
\begin{lstlisting}
# Limpieza y preparacion del dataset

df_prep = df.copy()

# 1) Verificacion de valores nulos
print(df_prep.isna().sum().sort_values(ascending=False))
print(f"Total de nulos: {df_prep.isna().sum().sum()}")

# 2) Verificacion de duplicados
print(f"Filas duplicadas completas: {df_prep.duplicated().sum()}")
print(f"Id duplicados: {df_prep['id'].duplicated().sum()}")

# 3) Eliminacion de columna identificadora para analisis/modelado
df_prep = df_prep.drop(columns=["id"])

# 4) Transformacion de la variable objetivo a formato numerico
df_prep["diagnosis_num"] = df_prep["diagnosis"].map({"B": 0, "M": 1})
\end{lstlisting}

\subsection{Justificación de cada decisión}
\begin{itemize}
    \item \textbf{Tratamiento de nulos:} no se imputaron valores porque, tras remover la columna vacía, no hubo nulos en las variables útiles.
    \item \textbf{Eliminación de duplicados:} no se eliminaron filas porque no se detectaron duplicados de filas ni de \texttt{id}.
    \item \textbf{Transformaciones:}
    \begin{itemize}
        \item Se eliminó \texttt{id} en el dataset de trabajo porque es un identificador, no una característica clínica.
        \item Se creó \texttt{diagnosis\_num} (B=0, M=1) para facilitar análisis cuantitativos y futuras etapas de modelado.
    \end{itemize}
\end{itemize}

\subsection{Preguntas guía}
\begin{itemize}
    \item \textbf{¿Qué pasa si NO limpian los datos?} Se pueden mantener errores de formato, columnas no informativas y estructuras inconsistentes que afecten análisis posteriores.
    \item \textbf{¿Qué decisiones podrían alterar los resultados?} Imputar nulos, eliminar filas/columnas, codificar variables o mantener identificadores puede cambiar el comportamiento de análisis y modelos.
    \item \textbf{¿Están perdiendo información al limpiar?} En este caso no se perdió información clínica relevante: se eliminó una columna vacía y una columna de identificación; se conservó el resto de variables clínicas.
\end{itemize}
